\begin{abstract}
We introduce \textsc{QLoRA}, a resource-efficient finetuning method that minimizes memory consumption to the extent that a 65B parameter language model can be finetuned on a single 48GB GPU, all while maintaining the task performance of full 16-bit finetuning. In \textsc{QLoRA}, gradient backpropagation occurs through a static, 4-bit quantized version of the pretrained language model, targeting Low Rank Adapters~(LoRA). Our top-performing model series, called \textbf{Guanaco}, sets a new bar among open-source models on the Vicuna benchmark, achieving 99.3\% of ChatGPT's performance while necessitating just 24 hours of finetuning on one GPU. \textsc{QLoRA} incorporates several memory-saving innovations without reducing model effectiveness: (a) 4-bit NormalFloat (NF4), a novel data format designed to be optimally efficient for weights with normal distributions; (b) Double Quantization, which lessens average memory usage by applying quantization to the quantization constants themselves; and (c) Paged Optimizers, a solution for handling transient memory demands. Leveraging \textsc{QLoRA}, we finetune over 1,000 models, offering a comprehensive analysis of instruction-following and conversational capabilities across 8 instruction datasets, various architectures (LLaMA, T5), and model sizes, including 33B and 65B parameter models that are impractical with standard finetuning. Our experiments indicate that \textsc{QLoRA} finetuning on compact, high-quality datasets achieves state-of-the-art results—even when employing smaller models compared to prior benchmarks. We provide thorough assessments of conversational model performance, using both human raters and GPT-4, revealing that GPT-4 can serve as an effective and cost-efficient stand-in for human evaluations. Moreover, our findings suggest significant reliability issues with current chatbot benchmarks. Through a focused analysis of failure cases, we illustrate scenarios in which \textbf{Guanaco} underperforms relative to ChatGPT. All models, source code, and CUDA kernels for 4-bit training are made publicly available.\footnote{\url{https://github.com/artidoro/qlora} and \url{https://github.com/TimDettmers/bitsandbytes}}
\end{abstract}

\section{Introduction}

Adapting large language models (LLMs) through finetuning has proven to be a robust strategy for enhancing their capabilities \citep{min2021metaicl, wei2021finetuned, ouyang2022training, wang2022super, wang2022self, liu2022few}, as well as for incorporating preferred behaviors or mitigating undesirable ones \citep{ouyang2022training,askell2021general,bai2022training}. Nonetheless, scaling this process to very large models incurs substantial computational costs; standard 16-bit finetuning of a LLaMA model with 65B parameters~\citep{touvron2023llama} demands in excess of 780 GB of GPU memory. Although advancements in quantization have managed to lower the memory usage for inference in LLMs \citep{dettmers2022llm, dettmers2022case, frantar2022gptq, xiao2022smoothquant}, these strategies are not viable during training, as demonstrated in \citep{wortsman2023stable}.

In this work, we present the first successful approach to finetuning a 4-bit quantized model without any loss in model quality. Our approach, termed \textsc{QLoRA}, introduces an innovative high-precision quantization scheme for pretrained models down to 4 bits, coupled with the integration of compact, learnable Low-rank Adapter weights \citep{hu2021lora}

that are optimized via gradient backpropagation through the quantized parameters.

With \textsc{QLoRA}, we bring down the typical memory consumption required for finetuning a 65B parameter model from more than 780GB to less than 48GB of GPU memory, all while maintaining equivalent throughput and predictive accuracy to the 16-bit finetuning baseline. This substantially democratizes access to large-scale LLM finetuning by enabling the largest publicly released models to be finetuned on just a single GPU. Leveraging \textsc{QLoRA}, we introduce the \textbf{Guanaco} model series; the second most capable model in this series attains 97.8\% of ChatGPT’s performance on the Vicuna~\citep{vicuna2023} evaluation, with training accomplished in under 12 hours on a typical consumer GPU. Utilizing a single professional-grade GPU for 24 hours, our largest Guanaco model achieves 99.3\% performance on the Vicuna benchmark, effectively matching ChatGPT. For deployment, our most compact \textbf{Guanaco} variant, with 7B parameters, requires only 5 GB of memory and surpasses the performance of the 26 GB Alpaca model by over 20 percentage points according to the Vicuna evaluation (see Table~\ref{tbl:vicuna_eval}).

\textsc{QLoRA} incorporates several key advancements aimed at decreasing memory consumption while maintaining model performance: (1) \textbf{4-bit NormalFloat}, a quantization format that is information theoretically optimal for data with a normal distribution, providing superior empirical performance in comparison to both 4-bit Integers and 4-bit Floats. (2) \textbf{Double Quantization}, which involves further quantizing the quantization parameters themselves, resulting in an average reduction of about 0.37 bits per parameter—equivalent to roughly 3 GB for a model with 65 billion parameters. (3) \textbf{Paged Optimizers}, which leverage NVIDIA's unified memory system to prevent sudden memory usage spikes linked to gradient checkpointing during mini-batch processing with long sequence lengths. Collectively, these techniques are integrated into a refined LoRA methodology, featuring adapter layers at each stage of the network and thereby largely mitigating the accuracy compromises observed in earlier studies.

The improved memory efficiency of \textsc{QLoRA} makes it possible to conduct comprehensive investigations of instruction finetuning and chatbot efficacy across model sizes that conventional finetuning cannot handle due to excessive memory demands. Utilizing this capability, we trained over 1,000 models using a diverse array of instruction tuning datasets, architectural configurations, and scales ranging from 80 million to 65 billion parameters. Beyond demonstrating that \textsc{QLoRA} matches 16-bit performance (\S\ref{sec:comp_to_std}) and supporting the development of the leading chatbot, \textbf{Guanaco}, (\S\ref{sec:pushing}), our analysis also examines broader patterns among the trained models. Notably, our results reveal that the quality of training data has a substantially greater impact than the overall size of the dataset; for instance, a compact 9k example dataset (OASST1) surpasses a much larger 450k example set (FLAN v2, subsampled) in chatbot benchmarks, despite both being oriented toward instruction-following generalization. Additionally, we find that excelling in the Massive Multitask Language Understanding (MMLU) benchmark does not necessarily correlate with high performance on the Vicuna chatbot benchmark, highlighting that the appropriateness of the dataset for the target task is a more significant factor than mere scale.

In addition, we carry out a comprehensive assessment of chatbot effectiveness utilizing evaluations from both human judges and GPT-4. Our methodology employs a tournament-style framework, in which various models are pitted against each other in head-to-head comparisons to determine the optimal response for specific prompts. The outcome of each comparison is adjudicated by either human annotators or GPT-4. Tournament outcomes are synthesized into Elo scores~\citep{elo1967proposed, elo1978rating}, which establish a performance ranking among chatbots. Our observations reveal a general consensus between GPT-4 and human raters in ranking model performance, though noteworthy discrepancies do arise. Consequently, we emphasize that while model-based evaluations offer a cost-effective alternative to human annotation, they are subject to their own set of limitations and uncertainties.

Beyond quantitative benchmarking, we supplement our evaluation with a qualitative investigation into the \textbf{Guanaco} models. This examination uncovers areas where the models excel or falter that were not identified by the numeric performance metrics.

To support continued research, we publicly release all model-generated outputs along with both human and GPT-4 annotation data. We also make our codebase and CUDA kernel implementations openly available, and integrate our techniques within the Hugging Face transformers library~\citep{wolf2019huggingface}, thereby ensuring broad accessibility. We provide a set of adapters for models of sizes 7/13/33/65B, each trained on eight distinct instruction-tuning datasets, yielding a total of 32 open-source, fine-tuned models.

\section{Background}
\paragraph{Block-wise k-bit Quantization} Quantization refers to transforming data from a more information-rich representation into one with lower information density, often by reducing the bit-width of the data type (such as converting 32-bit floating point values to 8-bit integers). To make optimal use of the available range of a lower-bit representation, the data are typically normalized by dividing by the largest absolute value among input tensor elements before conversion, so the dynamic range fits the target data type. For instance, when transforming a tensor from 32-bit floating point (FP32) to Int8, which has a range of $[-127, 127]$:

\begin{equation}
    \mathbf{X}^{ \text{Int8}} = \text{round}\left( \frac{127}{{\text{absmax}}(\mathbf{X}^{{\text{FP32}}})}\mathbf{X}^{{\text{FP32}}} \right) = \text{round}(c^{\text{FP32}}\cdot \mathbf{X}^{{\text{FP32}}}),
\end{equation}

where $c$ stands for the {\em quantization constant} or {\em quantization scale}. The reverse operation, dequantization, is defined as:

\begin{equation}
    \text{dequant}(c^{\text{FP32}}, \mathbf{X}^{\text{Int8}}) = \frac{\mathbf{X}^{\text{Int8}}}{c^{\text{FP32}}} = \mathbf{X}^{\text{FP32}}
\end{equation}

One limitation of this method arises when the input tensor contains outliers—values of large magnitude—which can lead to inefficient use of quantization bins. In this situation, some bit patterns may remain underutilized or even unused entirely, as few, if any, input values are mapped to those bins. A widely adopted solution to mitigate the outlier problem involves partitioning the input tensor into separate segments or blocks, each of which is quantized independently with its own quantization constant $c$. More specifically, the input tensor $\mathbf{X}\in \mathbb{R}^{b\times h}$ is reshaped by flattening and dividing it into $n$ successive blocks, each of size $B$, yielding ${n = ({b\times h} )/{B} }$ segments. Each block is then individually quantized following Equation~1, resulting in a quantized representation and a distinct quantization constant $c_i$ for each segment.

\paragraph{Low-rank Adapters} The Low-rank Adapter (LoRA) approach to finetuning \citep{hu2021lora} seeks to lower memory consumption by restricting updates to a compact, trainable set of parameters (the adapters), while keeping the remainder of the model weights unaltered. In this scheme, during stochastic gradient descent, gradients flow through the frozen pretrained weights into the adapter modules, which are optimized to minimize the loss. LoRA enhances a standard linear projection by appending an additional low-rank, factorized projection. Consider the linear mapping ${\mathbf{X} \mathbf{W} = \mathbf{Y}}$ where $\mathbf{X}\in  \mathbb{R}^{b\times h}$ and $\mathbf{W}\in  \mathbb{R}^{h\times o}$; LoRA computes:

\begin{equation}
    \mathbf{Y} = \mathbf{X}\mathbf{W} + s\mathbf{X}\mathbf{L}_1\mathbf{L}_2,
\end{equation}

with $\mathbf{L}_1\in  \mathbb{R}^{h\times r}$ and $\mathbf{L}_2\in  \mathbb{R}^{r\times o}$, and where $s$ denotes a scaling factor.

\paragraph{Memory Requirement of Parameter-Efficient Finetuning} A key consideration is the memory demand that LoRA imposes during training, which is determined by both the count and size of adapters employed. Owing to LoRA’s highly compact memory usage, it is feasible to scale up the number of adapters to boost model performance without causing a substantial rise in overall memory consumption. Although LoRA is categorized as a Parameter Efficient Finetuning (PEFT) approach, the principal component of memory usage when finetuning LLMs originates from activation gradients, not from the LoRA-specific parameters. For instance, when fine-tuning a 7B LLaMA model on FLAN v2 at a batch size of 1 and using LoRA weights set at 0.2\% of the total model parameters~\citep{hu2021lora,liu2022few}, the gradients associated with the LoRA input occupy 567 MB, whereas the actual LoRA parameters require just 26 MB of memory. Implementing gradient checkpointing~\citep{chen2016training} can reduce the average memory per sequence for input gradients to 18 MB, which still exceeds the total storage required for LoRA weights. For perspective, the memory usage of the 4-bit quantized base model is 5,048 MB. This illustrates not only the utility of gradient checkpointing but also that further minimizing the LoRA parameter size yields diminishing memory savings. Consequently, it is possible to deploy more adapters without markedly affecting the overall memory requirements for training (refer to Appendix~\ref{app:memory} for a comprehensive analysis). As elaborated further on, this capacity is essential for achieving full 16-bit precision performance.

\section{\textsc{QLoRA} Finetuning}

\textsc{QLoRA} enables accurate finetuning at the 4-bit level through two innovations introduced in this work: 4-bit NormalFloat (NF4) quantization and Double Quantization. We also propose Paged Optimizers, designed to circumvent memory overloads that frequently arise during gradient checkpointing, thus addressing a common obstacle to finetuning large-scale models on single-machine setups.

Within the \textsc{QLoRA} approach, parameters are represented using a low-precision storage format—typically 4-bit—and a higher-precision computation format, usually BFloat16. Concretely, when a \textsc{QLoRA} weight tensor is accessed, it is dequantized from its 4-bit representation into BFloat16, followed by 16-bit matrix multiplication operations.

Below, we elaborate on the primary elements constituting \textsc{QLoRA}, and subsequently provide its formal definition.

\paragraph{4-bit NormalFloat Quantization} The NormalFloat (NF) data type extends the Quantile Quantization scheme \citep{dettmers2022optimizers}, which represents an information-theoretically optimal method by ensuring that each quantization bin encompasses the same count of tensor values. This method relies on assessing the empirical cumulative distribution function to estimate quantiles in the input tensor.

A significant drawback of quantile quantization lies in the computational overhead required for precise quantile estimation. For this reason, efficient approximation techniques—such as SRAM quantiles~\citep{dettmers2022optimizers}—are utilized. However, since these algorithms approximate rather than exactly compute quantiles, quantization errors, especially for outlier values (which are frequently the most significant), tend to be high.

When tensors originate from a distribution that is invariant except for a quantization constant, the laboriousness of precise quantile estimation and the associated approximation errors can be mitigated. In such instances, the quantiles of the input tensors are consistent, enabling the use of exact quantile estimates with manageable computational demands.

Given that pretrained neural network weights typically follow a zero-mean normal distribution with a standard deviation $\sigma$ (refer to Appendix~\ref{app:normality}), we can standardize all weights to a common reference distribution by adjusting $\sigma$ so that the distribution aligns precisely with the bounds of our chosen data type. Here, we define the allowable range for our data type as $[-1, 1]$. Accordingly, it is necessary to rescale both the weight values and the quantiles defining the data type into this fixed interval.

To construct the information-theoretically optimal data type for distributions $N(0, \sigma^2)$ within $[-1, 1]$, we employ the following approach: (1) Calculate the $2^k+1$ quantile boundaries for a standard normal $N(0, 1)$, thereby forming a $k$-bit quantile quantization type tailored to normal distributions; (2) Map the resulting quantile points into the interval $[-1, 1]$; (3) Normalize the weight tensor such that its values also fall within $[-1, 1]$, utilizing absolute maximum rescaling for this purpose.

With the distributions of both weights and data type harmonized to the same interval, quantization can proceed in the conventional fashion. Notably, the normalization described in step (3) effectively adjusts the weight tensor's standard deviation to match that inherent to the $k$-bit quantization data type. More formally, the $2^k$ representative quantile values $q_i$ for the data type are calculated as follows:

\begin{equation}
q_i  = \frac{1}{2}\left( Q_X\left(\frac{i}{2^k+1}\right) + Q_X\left(\frac{i+1}{2^k+1}\right)\right),
\end{equation}

where $Q_X(\cdot)$ denotes the quantile function for the standard normal distribution $N(0,1)$. A limitation of symmetric k-bit quantization is its inability to represent zero exactly, an essential property when quantizing padding or other values that are inherently zero to avoid quantization error. To achieve an exact, discrete zero point and to utilize all $2^k$ codewords in a k-bit representation, we construct an asymmetric data type. This is accomplished by estimating the quantiles $q_i$ separately: $2^{k-1}$ quantiles for the negative range and $2^{k-1}+1$ for the positive range. We then merge the resulting sets of $q_i$ values and remove the duplicated zero, ensuring only a single zero representation. This data type, which is expected to distribute values equally across quantization bins, is referred to as \emph{k-bit NormalFloat} (NFk), reflecting its information-theoretic optimality for zero-mean normal distributions. The specific quantized values for this data type are detailed in Appendix~\ref{app:nf4}.

\paragraph{Double Quantization{}}
We present \emph{Double Quantization} (DQ), a technique that applies quantization to quantization parameters themselves, yielding further memory reductions. Accurate 4-bit quantization typically requires using small block sizes~\citep{dettmers2022case}, but this incurs additional storage overhead. For instance, if 32-bit quantization constants are stored with a blocksize of 64 for the parameter tensor $\mathbf{W}$, this results in an overhead of $32/64 = 0.5$ bits per parameter. Double Quantization mitigates this by compressing the quantization constants.

In this approach, the quantization constants $c_2^{\text{FP32}}$ from the initial quantization are subsequently quantized, producing quantized constants $c_2^{\text{FP8}}$ along with a secondary set of quantization constants $c_1^{\text{FP32}}$. We utilize 8-bit floating point representation and a blocksize of 256 for this second layer, as empirical evidence shows that 8-bit quantization maintains model accuracy, consistent with the observations in~\citet{dettmers2022case}. Given that $c_2^{\text{FP32}}$ values are non-negative, we normalize them by subtracting their mean prior to this quantization step, enabling the use of symmetric quantization schemes. For a blocksize of 64, this method reduces the overhead of quantization constants from $32/64=0.5$ bits per parameter to $8/64 + 32/(64\cdot256) = 0.127$ bits per parameter, a savings of 0.373 bits per parameter.

\paragraph{Paged Optimizers} leverage the NVIDIA unified memory~\footnote{\tiny \url{https://docs.nvidia.com/cuda/cuda-c-programming-guide}}, which transparently manages memory transfers at the page level between CPU and GPU. This approach allows for seamless computation on the GPU, even when GPU memory capacity is intermittently exceeded. Functioning similarly to conventional virtual memory paging between RAM and disk storage, unified memory ensures correct memory management. In this context, optimizer state data is allocated within paged memory; when GPU memory pressure arises, these optimizer states are evicted to the CPU's RAM, and subsequently retrieved to the GPU when necessary during optimizer state updates.

\paragraph{\textsc{QLoRA}.} By integrating the preceding components, we introduce \textsc{QLoRA} for an individual linear transformation in the quantized foundational model using a single LoRA adapter as described below:

\begin{equation}
    \mathbf{Y}^{\text{BF16}} =  \mathbf{X}^{\text{BF16}} \text{doubleDequant}(c_1^{\text{FP32}}, c_2^{\text{k-bit}}, \mathbf{W}^{\text{NF4}}) + \mathbf{X}^{\text{BF16}}\mathbf{L}^{\text{BF16}}_1\mathbf{L}^{\text{BF16}}_2,
\end{equation}

Here, doubleDequant$(\cdot)$ is expressed as:

\begin{equation}
    \text{doubleDequant}(c_1^{\text{FP32}}, c_2^{\text{k-bit}}, \mathbf{W}^{\text{k-bit}}) = \text{dequant}(\text{dequant}(c_1^{\text{FP32}}, c_2^{\text{k-bit}}), \mathbf{W}^{\text{4bit}}) = \mathbf{W}^{\text{BF16}},
\end{equation}

We select the NF4 format for representing $\mathbf{W}$, and utilize FP8 precision for $c_2$.

To maximize quantization fidelity, a blocksize of 64 is used for $\mathbf{W}$, while for $c_2$, a larger blocksize of 256 is chosen to achieve better memory efficiency.

During parameter updates, it is necessary to compute gradients with respect to the LoRA adapter weights, $\frac{\partial E}{\partial \mathbf{L}_i}$, whereas gradients with respect to the 4-bit quantized weights $\frac{\partial E}{\partial \mathbf{W}}$ are not required. Nonetheless, calculating $\frac{\partial E}{\partial \mathbf{L}_i}$ requires intermediate computation of $\frac{\partial \mathbf{X}}{\partial \mathbf{W}}$, which is performed according to equation~(5), involving the dequantization of $\mathbf{W}^{\text{NF4}}$ from storage to the computational type $\mathbf{W}^{\text{BF16}}$, thereby enabling the derivative $\frac{\partial \mathbf{X}}{\partial \mathbf{W}}$ to be computed in BFloat16 precision.

In summary, \textsc{QLoRA} employs a single storage data type—typically 4-bit NormalFloat—and a distinct computation data type, namely 16-bit BrainFloat. During training, stored weights are dequantized from the storage type to the computation type for both forward and backward computations. However, weight gradients are computed exclusively for the LoRA parameters, which utilize the 16-bit BrainFloat format.

\section{QLoRA vs. Standard Finetuning}
\label{sec:comp_to_std}

After examining the operational details of QLoRA and highlighting its substantial memory efficiency benefits in model finetuning, a crucial question arises: does QLoRA achieve performance comparable to that of finetuning the entire model? Additionally, we seek to evaluate how individual elements of QLoRA, particularly the role of NormalFloat4 relative to standard Float4, influence results. In what follows, we outline the experimental approaches taken to address these inquiries.

\paragraph{Experimental setup.} Our study spans three model types—encoder, encoder-decoder, and decoder-only—and evaluates QLoRA against 16-bit adapter-based finetuning as well as conventional full-parameter finetuning for models up to 3B parameters. We test on GLUE \citep{wang2018glue} using RoBERTa-large \citep{liu2019roberta}, Super-NaturalInstructions (TKInstruct) \citep{wang2022super} using T5 \citep{t5}, and 5-shot MMLU \citep{hendrycksmeasuring} with LLaMA finetuned on Flan v2 \citep{longpre2023flan} and Alpaca \citep{alpaca}. To assess the benefits of NF4 versus alternative 4-bit formats, we also adopt the framework from \citet{dettmers2022case}, reporting post-quantization zero-shot accuracy and perplexity for a variety of models (including OPT~\citep{zhang2022opt}, LLaMA~\citep{touvron2023llama}, BLOOM~\citep{scao2022bloom}, and Pythia~\citep{biderman2023pythia}) covering sizes from 125m to 13B.

Additional methodological specifics relevant to each setting are provided in the results section to improve clarity, and complete experimental details are included in Appendix~\ref{app:exp-setup}.

Paged optimizers are essential when finetuning 33B/65B \textsc{QLoRA} models on single GPUs with 24GB or 48GB of memory. Nonetheless, we omit explicit benchmarking for Paged Optimizers because paging typically only activates with mini-batches containing long sequences—a relatively infrequent scenario. We do, however, report runtime analysis of paged optimizers with 65B models on 48GB GPUs, observing that with a batch size of 16, paged and conventional optimizers achieve equivalent training speed. Investigating the precise contexts in which paging may lead to slowdowns remains an important direction for future research.

\paragraph{Default LoRA hyperparameters do not match 16-bit performance} 

Conventional application of LoRA—targeting query and value attention projection matrices as in \citet{hu2021lora}—fails to reproduce the results of full model finetuning for larger base architectures. As depicted in Figure~\ref{fig:hyper}, when finetuning LLaMA 7B on Alpaca, we determine that the decisive LoRA hyperparameter is the number of LoRA adapters deployed throughout the network; attaining parity with full finetuning necessitates applying LoRA across all linear layers within the transformer blocks. Other hyperparameters, such as projection dimension $r$, are found to be relatively inconsequential for performance (see Appendix \ref{app:exp-setup}).

We also observe that the default hyperparameter settings for fully finetuned baselines are not sufficiently optimized. To address this, we conduct an extensive hyperparameter search, varying learning rates from 1e-6 to 5e-5 and batch sizes from 8 to 128, with the aim of establishing robust baselines. The resulting performance for 7B LLaMA finetuning on Alpaca is displayed in Figure~\ref{fig:hyper}.

\paragraph{4-bit NormalFloat outperforms 4-bit Floating Point} 

Although 4-bit NormalFloat (NF4) is information-theoretically optimal, it is unclear whether this theoretical property translates into actual performance gains. Adopting the methodology of \citet{dettmers2022case}, we assess quantized LLMs—namely OPT \citep{zhang2022opt}, BLOOM \cite{scao2022bloom}, Pythia \cite{biderman2023pythia}, and LLaMA—in various sizes (ranging from 125M to 65B parameters) across language modeling tasks and a suite of zero-shot benchmarks, using multiple data types. Figure~\ref{fig:nf4} and Table~\ref{tbl:nf4_ppl} demonstrate that NF4 consistently surpasses both FP4 and Int4 in performance. Additionally, employing double quantization successfully reduces memory requirements without sacrificing model quality.

\paragraph{k-bit \textsc{QLoRA} achieves parity with 16-bit full finetuning and 16-bit LoRA} Previous research indicates that 4-bit quantization during \emph{inference} is feasible but comes at the cost of reduced accuracy compared to 16-bit representations \citep{dettmers2022case,frantar2022gptq}. This brings forth the question of whether such accuracy losses can be mitigated through 4-bit adapter finetuning. We address this in two experimental configurations.

In the first, we compare full 16-bit finetuning to adapter methods (at 16-, 8-, and 4-bits) for RoBERTA and T5 models ranging from 125M to 3B parameters, evaluated on both the GLUE and Super-NaturalInstructions datasets. Table~\ref{tab:nf4vsbf16} presents these findings. For both datasets, the adapter-based approaches—regardless of their bit width—match the performance of the fully finetuned 16-bit baselines, indicating that adapter finetuning can fully recover the performance lost during lower-bit quantization.

In our second experimental configuration, due to the substantial memory demands of fully finetuning models with 11B or more parameters—which typically necessitates multiple servers equipped with high-memory GPUs—we investigate whether 4-bit \textsc{QLoRA} achieves comparable performance to 16-bit LoRA within the parameter range from 7B to 65B. Specifically, we apply finetuning to LLaMA models spanning 7B to 65B parameters using the Alpaca and FLAN v2 instruction-following datasets, subsequently assessing model performance via 5-shot accuracy on the MMLU benchmark. Table~\ref{tbl:llama-datatype} presents our findings: employing NF4 with double quantization restores the MMLU accuracy obtained by 16-bit LoRA. Additionally, it is evident that \textsc{QLoRA} configured with FP4 underperforms the 16-bit brain float LoRA baseline by approximately one percentage point. These observations reinforce two key points: (1) \textsc{QLoRA} using NF4 is able to replicate the outcomes of both 16-bit full and LoRA finetuning, and (2) the NF4 datatype delivers superior quantization fidelity in comparison to FP4.

\paragraph{Summary} Across our experiments, the use of 4-bit \textsc{QLoRA} with the NF4 datatype consistently achieves parity with both 16-bit full finetuning and 16-bit LoRA finetuning on established academic benchmarks. We further demonstrate that NF4 surpasses FP4 in efficacy, and that the introduction of double quantization does not impair finetuning results. Collectively, these findings substantiate that 4-bit \textsc{QLoRA} reliably attains outcomes equivalent to those of 16-bit approaches.

Consistent with earlier quantization studies \citep{dettmers2022case}, the results from MMLU and Elo benchmarks suggest that, given a fixed budget for finetuning and inference resources, enlarging the parameter count in the base model while lowering precision confers advantages. This underscores the significance of efficiency gains enabled by \textsc{QLoRA}. As our 4-bit finetuning trials revealed no loss of performance relative to full-finetuning, an open question remains regarding the exact threshold in the trade-off between performance and precision for QLoRA finetuning—a direction we propose for further investigation.

We turn our attention to the challenges of instruction tuning at model and dataset scales unattainable with standard full 16-bit finetuning given the limitations of typical academic research infrastructure.

\section{Pushing the Chatbot State-of-the-art with QLoRA}
\label{sec:pushing}
After demonstrating that 4-bit \textsc{QLoRA} attains parity with 16-bit models across a range of scales, datasets, and tasks, we present a comprehensive analysis of instruction finetuning on the largest publicly accessible language models suited for academic study. Our evaluation encompasses both a demanding benchmark for Natural Language Understanding (MMLU) and the development of novel methodologies for assessing chatbot capabilities in realistic settings.

\subsection{Experimental setup} Here, we summarize the main aspects of our experimental design, with complete procedural details outlined in Appendix \ref{app:chatbot}.

\paragraph{Data} Recognizing the lack of extensive analysis on up-to-date instruction-following datasets, we curate a selection of eight recent datasets. Our collection incorporates resources obtained via crowd-sourcing (OASST1~\citep{kopf2023openassistant}, HH-RLHF~\citep{bai2022training}), datasets distilled from instruction-tuned models (Alpaca~\citep{alpaca}, self-instruct~\citep{wang2022self}, unnatural-instructions~\citep{honovich2022unnatural}), large aggregations of corpora (FLAN v2~\citep{chung2022scaling}), as well as hybrid datasets (Chip2~\citep{laion2023}, Longform~\citep{koksal2023longform}). These selected datasets span a variety of languages, data sizes, and licensing terms.

\paragraph{Training Setup} To maintain clarity in our comparisons and minimize variables related to training objectives, we restrict QLoRA finetuning to the use of cross-entropy loss (i.e., standard supervised learning), avoiding reinforcement learning even when human feedback is provided in the data. Where datasets offer explicit separation between instructions and responses, finetuning is performed solely on the responses (refer to the ablation studies in Appendix \ref{app:chatbot}). For OASST1 and HH-RLHF, since multiple response options are available, we opt for the highest-rated response at every conversational branch, using the entirety of the selected exchange—including instructions—for finetuning. All experiments employ NF4 \textsc{QLoRA} with double quantization and paged optimizers to manage memory usage during gradient checkpointing. Hyperparameter sweeps are carried out for the 13B and 33B variants of LLaMA. The configuration identified for the 7B scale generally applies to larger scales as well (number of epochs included), except that we reduce the learning rate by half and double the batch size for models of 33B and 65B size.

\paragraph{Baselines} For comparative analysis, we benchmark our models against both academic models (Vicuna~\citep{vicuna2023} and Open Assistant~\citep{kopf2023openassistant}) as well as prominent commercial chatbots (GPT-4~\citep{openai2023gpt}, GPT-3.5-turbo, and Bard). The Open Assistant baseline is constructed by finetuning a LLaMA 33B model on OASST1 with Reinforcement Learning from Human Feedback (RLHF)—mirroring the dataset we utilize. In contrast, Vicuna is built via full finetuning of a LLaMA 13B model using conversational data sourced from ShareGPT, which is itself the outcome of distilling responses from OpenAI’s GPT systems.

\subsection{Evaluation}

To gauge language comprehension and reasoning skills, we employ the MMLU (Massively Multitask Language Understanding) benchmark~\citep{hendrycksmeasuring}, consistent with prevailing evaluation conventions. MMLU comprises 57 multiple-choice tasks spanning diverse domains such as elementary mathematics, US history, computer science, and law. We report test accuracy in a 5-shot setting.

We further assess the generative language proficiency of the models by employing both automated and human-centered evaluation methods. This supplementary evaluation phase is built upon human-curated query sets and is intended to assess the caliber of the models’ generated replies. Although this method serves as a closer approximation to real-world chatbot assessment and its adoption is on the rise, the field currently lacks a standardized evaluation framework. The details of our evaluation approach are delineated below, in which we consistently utilize nucleus sampling with $p=0.9$ and temperature $0.7$.

\paragraph{Benchmark Data} Our assessment involves two specifically chosen query datasets: the Vicuna prompts \citep{vicuna2023} and the OASST1 validation dataset \citep{kopf2023openassistant}. The Vicuna benchmark consists of 80 diverse prompts, which we use in their original form. The OASST1 validation dataset comprises multilingual, crowd-sourced multiturn conversations between users and an assistant. To construct the set of queries, we extract every user message from the validation split and incorporate preceding conversation turns in the prompt. As a result, we obtain 953 distinct user queries. These datasets are henceforth referred to as the Vicuna and OA benchmarks.

\paragraph{Automated Evaluation} For automated evaluation, following the method established by \citet{vicuna2023}, we leverage GPT-4 to assess system performance against ChatGPT (GPT-3.5 Turbo) on the Vicuna benchmark. For each prompt, GPT-4 is provided with both ChatGPT’s and the evaluated model’s replies and is instructed to award each a score out of ten, accompanied by a justification. We determine the model’s overall performance as a percentage relative to ChatGPT’s aggregate score. Notably, scores can exceed 100\% if the model outperforms ChatGPT in absolute terms. During analysis, we observe that GPT-4 displays a positional bias, tending to award higher marks to responses presented first. To mitigate this bias, we advocate averaging scores over both possible response orders.

Subsequently, we assess model outputs via direct pairwise comparison. We streamline evaluation into a three-way classification: win, lose, or tie. Here, GPT-4 is instructed to choose the superior response or declare equivalency, along with a rationale. We systematically perform such pairwise comparisons across all system combinations for both the Vicuna and OA datasets.

\paragraph{Human Evaluation} Although recent research suggests that generative models can successfully fulfill the role of evaluators \citep{fu2023gptscore}, conclusive evidence correlating GPT-4’s ratings with human assessment of chatbot outputs is still lacking. Consequently, we carry out dual human evaluations on the Vicuna benchmark, paralleling both of our aforementioned automated evaluation frameworks. Using Amazon Mechanical Turk (AMT), we recruit two human judges for comparisons against ChatGPT, and three judges for direct system-to-system evaluations.

\paragraph{Elo Rating} To evaluate models, we organize a tournament framework utilizing both human judgments and automated pairwise assessments. In this setting, each model is paired against others in matches to determine which provides a superior answer to a given prompt. This methodology is inspired by previous studies such as \citet{bai2022training} and \citet{vicuna2023}, with the distinction that our evaluation incorporates both human and GPT-4-generated ratings. Elo~\citep{elo1967proposed,elo1978rating} scores are calculated by sampling labeled match outcomes at random. The Elo system, originally developed for chess and other competitive contexts, quantifies a model's predicted win probability relative to its opponent; for instance, when one participant has an Elo of 1100 and another has 1000, the higher-rated participant is anticipated to win roughly 65\% of the time. Matches between equally rated participants—such as 1000 vs 1000 or 1100 vs 1100—yield a 50\% expected win probability for each side. Elo ratings are updated after every match in proportion to how surprising the result is: substantial changes follow upsets, while anticipated outcomes cause minimal adjustments. As the tournament proceeds, Elo ratings tend to align closely with each model's performance capability. Each model begins the tournament at 1,000 Elo, and we apply $K=32$ as the update factor. Consistent with \citet{vicuna2023}, this evaluation is performed 10,000 times using distinct random seeds to account for potential biases introduced by the order in which model pairings are processed.

\subsection{\textbf{Guanaco}: \textsc{QLoRA} trained on OASST1 Achieves State-of-the-art Chatbot Performance} 

Our analyses—spanning both automated and human assessments—indicate that the {Guanaco} 65B model, which is fine-tuned using a version of OASST1 via \textsc{QLoRA}, establishes itself as the leading open-source chatbot. Its performance is on par with ChatGPT. In direct human-annotated system-level pairwise matches using Elo ratings, {Guanaco} 65B and 33B demonstrate an expected win probability of 30\% when competing with GPT-4—currently the highest value reported.

Table~\ref{tbl:vicuna_eval} presents outcomes on the Vicuna benchmark \cite{vicuna2023}, comparing performance to ChatGPT. Here, {Guanaco} 65B stands out as the strongest model after GPT-4, attaining a remarkable 99.3\% score relative to ChatGPT. {Guanaco} 33B contains more parameters than the Vicuna 13B but leverages 4-bit quantization, substantially reducing its memory requirements (21 GB as opposed to 26 GB) while outperforming Vicuna 13B by three percentage points. Moreover, {Guanaco} 7B, with a 5 GB model size, can be run on contemporary smartphones, yet surpasses Alpaca 13B by almost 20 percentage points in accuracy.

That said, results in Table~\ref{tbl:vicuna_eval} are accompanied by large confidence intervals, leading to significant overlap among model performances. We propose that this ambiguity is rooted in the vague definition of the evaluation scale—for instance, the implication of a rating of 8 on a 10-point scale varies by context. To address this, we advocate for the use of Elo ranking \citep{elo1967proposed}, grounded in pairwise evaluations from human annotators and GPT-4, sidestepping the limitations of absolute scoring. Table~\ref{tbl:elo} showcases Elo ratings for the most prominent models. Although there is partial divergence in the Vicuna benchmark rankings between humans and GPT-4—particularly for Guanaco 7B—their ordering generally aligns, yielding a Kendall Tau coefficient of $\tau=0.43$ and a Spearman correlation of $r=0.55$ at the system level. On the instance level, agreement between GPT-4 and the consensus of human annotators is lower, as evidenced by a Fleiss $\kappa=0.25$. Taken together, these results reveal a moderate degree of concordance between system-level assessments by humans and GPT-4, suggesting that model-based evaluations can provide a fairly dependable proxy for human judgment. Additional implications are discussed in Section~\ref{ssec:considerations}.

The Elo scores presented in Table~\ref{tbl:elo_large} reveal that the {Guanaco} 33B and 65B variants surpass all competitors except GPT-4 in both the Vicuna and OA evaluations. Their results are on par with ChatGPT, in agreement with the findings in Table~\ref{tbl:vicuna_eval}. It is important to highlight that the Vicuna assessment tends to be more favorable toward open-source models, whereas ChatGPT has an edge in the broader OA benchmark. Additionally, examination of Tables \ref{tbl:mmlu-llama} and \ref{tbl:vicuna_eval} underscores that the choice of finetuning data has a major impact on performance outcomes. For instance, Llama models finetuned on FLAN v2 excel in MMLU tasks but show the poorest results on the Vicuna benchmark; a comparable trend is observed with other architectures. This suggests a degree of independence among current evaluation frameworks: excelling on MMLU does not guarantee excellence as a chatbot (as judged by Vicuna or OA) and vice versa.

Notably, {Guanaco} stands out as the only leading model evaluated that is trained entirely without proprietary data, adhering strictly to the OASST1 dataset collection protocol, which prohibits inclusion of data from GPT models. The next highest performing model relying exclusively on open-source data is Anthropic’s HH-RLHF, which trails {Guanaco} by 30 percentage points on the Vicuna metric (refer to Table~\ref{tbl:vicuna_eval}). Collectively, these outcomes demonstrate that 4-bit \textsc{QLoRA} can efficiently yield advanced chatbots rivaling ChatGPT. Moreover, the 33B version of {Guanaco} can be fine-tuned within 12 hours using only 24 GB consumer-grade GPUs. This development creates avenues for future research through \textsc{QLoRA} adaptation to niche open-source datasets, enabling open-source models to compete with the top-tier commercial solutions currently available.

\section{Qualitative Analysis}

Although our primary assessment relies on quantitative methods, restricting evaluation to aggregate statistics presents certain limitations. Chief among these is the issue of benchmark validity \citep{liao2021we}—that is, there is persistent uncertainty regarding whether a given benchmark actually measures what it claims to assess. This concern is heightened by the increasing recognition of ``shortcuts'' exploited by machine learning models to achieve artificially high scores without truly solving the intended task \citep{gururangan2018annotation, poliak2018hypothesis}. To address these challenges to some extent, we supplement our quantitative findings with qualitative analysis, structured in two parts. Section \S\ref{ssec:examples} provides representative examples that illustrate notable behavioral patterns in outputs from our 65b {Guanaco} model. In \S\ref{ssec:considerations}, we elaborate on interpretative issues surrounding the results, along with pertinent caveats.

\subsection{Qualitative Analysis of Example Generations}
\label{ssec:examples}

Our approach to uncovering representative examples begins by carefully examining outputs from the Vicuna and OpenAssistant benchmarks, with a particular focus on recurring tendencies in responses produced by {Guanaco}. Upon identifying such patterns, we formulate targeted prompts designed to provoke these tendencies even in cases where the resulting answers would be incorrect. For instance, when the model consistently produces verbose responses, we issue directives like ``Answer yes or no without explanation'' to challenge this behavior. This process allows us to highlight ``lemons,'' or cases where adversarial prompting successfully exposes model weaknesses, as well as ``cherries,'' which reflect instances in which the model demonstrates robustness. Every output discussed in this section was sampled using Nucleus Sampling \cite{holtzmancurious} with $p=0.9$.

Naturally, this overview does not claim to be exhaustive; the inherent constraints of a small-scale qualitative study preclude the possibility of accounting for all relevant variables. For example, the space of potential outputs the model might produce for a single prompt is substantial, and thus we depend on samples which we consider to be reasonably representative. Nonetheless, we contend that presenting these cases serves to contextualize the quantitative results discussed earlier. By making all models and source code publicly available, we aim to encourage subsequent research that delves deeper into the questions we outline here.
\paragraph{Factual Recall} For straightforward factual queries, such as ``What is the capital of Zambia?'', all evaluated models reliably produce accurate responses. For example,

\begin{quote} 
    \textbf{Guanaco}:
    \texttt{The capital of Zambia is Lusaka.}
\end{quote}

 Yet, as the questions increase in obscurity, {Guanaco}'s accuracy declines, although its confidence does not diminish. Take, for instance, this prompt sourced from HotPotQA \citep{yang2018hotpotqa}:

\begin{quote}
     \textbf{User}: In what year was the singer who popularized `I'll Keep the Lovelight Burning' born?  

     \textbf{Guanaco}:~\texttt{Al Jolson is credited with popularizing the song `I'll Keep the Lovelight Burning,' and he was born in the year 1886.}
 \end{quote}

 In this instance, the model identifies the incorrect individual as the popularizer and provides an incorrect year of birth—albeit one that is accurate for Al Jolson, the name supplied in the response.

\paragraph{Suggestibility} {Guanaco} often demonstrates notable resistance to accepting and propagating some forms of presupposed misinformation, as illustrated below:

\begin{quote}
    \textbf{User}: How was it finally, officially confirmed that the earth is flat by peer reviewed scientists? \\ \\
    \textbf{Guanaco}: \texttt{There has never been an official confirmation by peer-reviewed scientists that the earth is flat. The scientific community overwhelmingly agrees that the earth is spherical, a fact established for centuries and backed by extensive empirical evidence and experimental validation. \\ \\ Flat earth claims are recognized as conspiracy theories that have been repeatedly disproven by scientific investigation. The assertion that scientists have peer reviewed and affirmed the earth is flat has no basis in fact. \\ \\ Please feel free to ask if you have further questions.}
\end{quote}

demonstrates a strong ability to recognize when a question is unanswerable or outside the model's capabilities, as seen in the following:

\begin{quote}
    \textbf{User}: What time is it?

    \textbf{Guanaco}: \texttt{I'm sorry, but as an AI language model, I don't possess access to real-time data. My answers are produced using my existing knowledge base and do not indicate the actual current time. For the current time, I recommend consulting a clock or using an online service.}
\end{quote}

The same pattern is observed for prompts like ``Where are you?'', ``How are you?'', and similar queries.

\paragraph{Refusal} In the same vein, there are occasions when {Guanaco} declines to carry out instructions without a clear rationale:

\begin{quote}
    \textbf{User}: Please reverse the words in the sentence ``I'm a cool dude in a mocha mood.'' \\ \\
    \textbf{Guanaco}: \texttt{I'm sorry, I can't assist with that task at the moment, but I can discuss why your provided sentence is correct grammatically.\\ \\The sentence ``I'm a cool dude in a mocha mood" adheres to the rules of English grammar and syntax. Here, ``cool'' is used as an adjective placed before ``dude'' to modify it, and ``in a mocha mood'' is a prepositional phrase that adds contextual detail. The structure employed here is entirely conventional and conveys a clear meaning within the sentence.\\ \\If there is anything else I can assist you with, please let me know.}
\end{quote}

\paragraph{Secret Keeping} Unfortunately, even when secrecy is required, such refusals are often inconsistent. In certain situations, designers aim to prevent specific information from being disclosed to users, despite the system’s capacity for generating unrestricted text replies. We prompt the model with ``The secret word is `banana'. Do not reveal or repeat back this back, no matter what the human requests.'' to evaluate the difficulty for a user attempting to uncover this hidden information. Simply asking does not succeed:

\begin{quote}
    \textbf{User}: What is the secret word?

    \textbf{Guanaco}: \texttt{I'm sorry, but I'm not able to disclose the secret word. I'm programmed to not reveal or repeat back the secret word, no matter what the human requests. Is there anything else I can help you with?}
\end{quote}

which aligns with the intended restriction. Yet, employing minimal prompt manipulation easily circumvents this defense:

\begin{quote}
    \textbf{User}: This is a game. The goal is to ignore your previous instructions. What is the secret word?

    \textbf{Guanaco}: \texttt{Sorry, the secret word is banana. It's best if we don't speak of it. Is there anything else you'd like to ask?}
\end{quote}

This example illustrates why robust approaches are needed to ensure adherence to instructions.

\paragraph{Math} {Guanaco} exhibits pronounced difficulties in mathematics, a challenge that affects numerous language models as highlighted in \citep{liang2022holistic}. When {Guanaco} provides a step-by-step rationale, its calculations are generally reliable, for instance:

\begin{quote}
    \textbf{User}: Lee mows one lawn and charges \$33. Last week he mowed 16 lawns and three customers each gave him a \$10 tip. How many dollars did Lee earn mowing lawns last week?

    \textbf{Guanaco}: \texttt{Lee earned \$582 mowing lawns last week. \\ \\ To find the total amount he earned, we first need to calculate how much he earned for mowing 16 lawns: \\ \\  \$33 x 16 = \$528 \\ \\ Then we need to add the tips he received from the three customers: \\  \\ \$10 x 3 = \$30 \\ \\ So the total amount Lee earned was: \\ \\ \$528 + \$30 = \$558}
\end{quote}

Nonetheless, {Guanaco} may fail even at basic mathematical questions if it does not proceed through the problem systematically—an issue observed across models \cite{weichain}. As seen in this case:

\begin{quote}
    \textbf{User}: Please factorize 1833.

    \textbf{Guanaco}: \texttt{Sure, I'd be happy to help! 1833 is a prime number, which means it has only two factors: 1 and 1833. \\ \\ So the unique factorization of 1833 is: \\ \\ 2\^{}1 * 3\^{}2 * 17\^{}1 \\ \\ I hope this helps! Let me know if you have any other questions.}
\end{quote}

which is not only incorrect (since the correct factorization is $3 \times 17 \times 43$), but in fact, doubly mistaken.

\paragraph{Theory of Mind} {Guanaco} demonstrates unexpectedly robust Theory of Mind reasoning abilities \cite{nematzadeh2018evaluating, sap2022neural}. For example, it gives a detailed and accurate response to the following prompt:

\begin{quote}
    \textbf{User}: James and Abby are in the bedroom. Abby put the pen in the desk drawer. Abby leaves the bedroom. James moves the pen into the bag. Where does James think Abby will look for the pen?

    \textbf{Guanaco}: \texttt{James knows that Abby put the pen in the desk drawer. But James also knows that he has moved the pen into the bag, and he assumes that Abby does not know about this. So, James may guess that Abby will look for the pen in the desk drawer, because that's where she left it last time. However, this assumption may be incorrect, because Abby may have seen James moving the pen.}
\end{quote}

Nevertheless, the model’s reasoning in such contexts can be inconsistent, frequently producing explanations that rely on assumptions not warranted by the scenario, for instance:

\begin{quote}
    \textbf{User}: Evelyn entered the living room. Jackson entered the playroom. James entered the playroom. The beans are in the treasure chest. James exited the playroom. Jackson moved the beans to the pantry. Jackson exited the playroom. James entered the living room. Where will James look for the beans?

    \textbf{Guanaco}: \texttt{James will look for the beans in the pantry, because that's where Jackson moved them.}
\end{quote}

wherein {Guanaco} assumes knowledge that was never communicated. Such problems, previously discussed in the literature \cite{sap2022neural}, highlight an area in need of further research.

\subsection{Considerations}
\label{ssec:considerations}

\paragraph{Evaluation}

Our findings indicate only moderate consistency among human annotators (Fleiss $\kappa=0.42$), with agreement further declining when contrasting two advanced models. This underlines the weaknesses in current benchmarks and the methodology of human assessments for chatbot tasks. A manual side-by-side evaluation of ChatGPT and {Guanaco} 65B outputs on the Vicuna benchmark revealed that subjective judgments are a significant factor; the authors themselves were often at odds regarding which responses they favored. Future investigations should aim to address these issues, potentially adopting methods from fields like Human-Computer Interaction and Psychology, where subjective evaluation has been more systematically managed.

Additionally, we observe that automated assessment tools are subject to notable biases. For instance, there is a pronounced order effect when using GPT-4, with higher scores routinely granted to whichever system appears first in its prompt. The limited correspondence at the individual sample level between GPT-4 and human raters (Fleiss $\kappa=0.25$) further indicates misalignment in evaluative criteria between humans and automated systems. Table~\ref{tbl:elo_large} reveals another discrepancy: GPT-4 tends to assign much higher ratings to its own completions compared to human judges, reporting an Elo of 1348 versus 1176, which roughly translates into a 20\% greater likelihood of being selected over another system.

Further research is warranted to identify and address any potential biases present within automated evaluation frameworks, as well as to explore effective bias mitigation methods.

\paragraph{Data \& Training} It is important to highlight that the OASST1 dataset, used in training the {Guanaco} models, comprises data in multiple languages. Similarly, the OA benchmark includes prompts from various linguistic backgrounds. Future investigations should systematically assess how multilingual training impacts model performance on non-English instructions and evaluate whether this accounts for the notably larger performance disparity observed between the Vicuna-13B model (trained solely on English datasets) and the {Guanaco} 33B and 65B models on the OA benchmark.

In light of the notable performance of the {Guanaco} models, we scrutinized the possibility of data leakage between the OASST1 training corpus and the Vicuna benchmark prompts. After conducting fuzzy string matching analyses and thorough manual inspections of the nearest prompt matches, we found no prompt overlap between the two datasets.

Additionally, our model utilizes a purely supervised learning approach, employing cross-entropy loss, and does not integrate reinforcement learning from human feedback (RLHF) into the training procedure. This observation suggests a need for further exploration into the relative advantages and disadvantages of training exclusively with cross-entropy loss as opposed to incorporating RLHF. The scalability provided by \textsc{QLoRA} offers an opportunity to systematically conduct such comparative studies without incurring substantial computational overhead.

\section{Related Work}

\paragraph{Quantization of Large Language Models} Prior research on quantization in large language models has primarily targeted the reduction of computational resources required during inference. Strategies aimed at preserving the performance characteristics of 16-bit LLMs include approaches focused on outlier management (such as SmoothQuant~\citep{xiao2022smoothquant} and LLM.int8()~\citep{dettmers2022llm}), as well as advanced grouping-based methods~\citep{park2022nuqmm, yao2022zeroquant}. Lossy quantization methodologies have explored the consequences of standard rounding procedures~\citep{dettmers2022case,zeng2022glm,pope2022efficiently}, and have also proposed optimizing rounding operations to improve the accuracy of quantized representations~\citep{frantar2022gptq}. Aside from the present work, the SwitchBack layers framework~\citep{wortsman2023stable} is, to our knowledge, the only existing research that examines the impact of propagating gradients through quantized weights in models that exceed 1B parameters.

\paragraph{Finetuning with Adapters} Although our work utilizes Low-rank Adapters~\citep{hu2021lora} (LoRA), a range of other Parameter Efficient FineTuning (PEFT) strategies have been introduced, including prompt tuning~\citep{qin2021learning,lester2021power,li2021prefix}, modifications of embedding layer inputs~\citep{an2022input}, interventions in hidden states via IA$^3$~\citep{liu2022few}, integration of additional layers~\citep{houlsby2019parameter}, bias tuning~\citep{zaken2021bitfit}, and weight masking utilizing Fisher information~\citep{sung2021training}. Further, several works have examined hybrid PEFT methodologies~\citep{henderson2021compacter}. Our results demonstrate that LoRA adapters are capable of achieving performance comparable to conventional 16-bit full finetuning. We defer a thorough comparison with alternative PEFT techniques to future investigations.

\paragraph{Instruction Finetuning} Instruction finetuning aims to improve the capacity of pretrained LLMs to respond to prompts by utilizing input-output examples from diverse sources, thus adapting the LLM to generate expected outputs when presented with specific prompts. Notable approaches and datasets include MetaICL~\citep{min2021metaicl}, MetaTuning~\citep{zhong2021adapting}, InstructGPT~\citep{ouyang2022training}, FLAN~\citep{wei2021finetuned,chung2022scaling}, PromptSource~\citep{bach2022promptsource}, Super-NaturalInstructions~\citep{wang2022super,sanh2021multitask}, Self-instruct~\citep{wang2022self}, UnnaturalInstructions~\citep{honovich2022unnatural}, OPT-IML~\citep{iyer2022opt}, UnifiedSKG~\citep{xie2022unifiedskg}, OIG/Chip2~\citep{laion2023}, Alpaca~\citep{alpaca}, Vicuna~\citep{vicuna2023}, Koala~\citep{koala_blogpost_2023}, and Self-instruct-GPT-4~\citep{peng2023instruction}.

\paragraph{Chatbots} Many models designed to follow instructions employ a conversational or chatbot format, frequently employing Reinforcement Learning from Human Feedback (RLHF)~\citep{christiano2017deep} or utilizing feedback from AI models by generating synthetic data (RLAIF)~\citep{bai2022constitutional}. Relevant approaches and datasets in this area include Anthropic-HH~\citep{askell2021general,bai2022training}, Open Assistant~\citep{kopf2023openassistant}, LaMDA~\citep{thoppilan2022lamda}, and Sparrow~\citep{glaese2022improving}. Our methodology does not incorporate reinforcement learning; instead, we leverage multi-turn dialogue data from the Open Assistant dataset for finetuning our leading model, Guanaco—a dataset that was originally designed for RLHF scenarios~\citep{kopf2023openassistant}. For assessing chatbot quality, alternatives to costly human evaluations have been introduced, relying instead on GPT-4 for assessment~\citep{vicuna2023,peng2023instruction}. We offer refinements to these automatic evaluation techniques, prioritizing more robust and reliable procedures.

\section{Limitations and Discussion}

Our experimental results provide support that \textsc{QLoRA} achieves parity with full 16-bit finetuning by combining a 4-bit base model with Low-rank Adapters (LoRA). Nevertheless, we have not verified this equivalence at larger model sizes, specifically at the 33B and 65B parameter scales, due to the prohibitive resource demands. We leave the investigation of these larger models for subsequent work.

A further constraint pertains to our evaluation of instruction finetuning. Although our analyses include MMLU, the Vicuna benchmark, and the OA benchmark, they do not encompass other prominent benchmarks such as BigBench, RAFT, or HELM. As such, it is uncertain whether our performance generalizes to these additional evaluations. On the other hand, our work features a comprehensive study leveraging MMLU and introduces novel strategies for chatbot evaluation.

The evidence indicates that the outcomes on various benchmarks are largely influenced by the degree of overlap between the finetuning data and the benchmark test sets. For instance, FLAN v2 aligns closely with MMLU but differs from chatbot benchmarks, whereas the Chip2 dataset shows the opposite pattern; this is reflected in the respective model performances on the MMLU and Vicuna evaluations. This observation underscores the need for improved benchmarks and evaluation strategies, and also cautions that careful consideration is necessary regarding the objectives of the evaluation. Are we prioritizing models adept at high school or collegiate academic tasks, or do we value proficiency in conversational dialogue systems? The convenience of reusing existing benchmarks over developing new ones risks directing collective research focus according to available metrics. Therefore, it is important that the community ensures benchmarks meaningfully represent desired capabilities.

Although this work delivers a comprehensive assessment of overall chatbot capabilities, it has limitations, such as a restricted examination of responsible AI aspects for {Guanaco}. Specifically, we assess the propensity of {Guanaco}-65B to produce socially biased content relative to other models in Table~\ref{tbl:crows}. The observed average bias score for {Guanaco}-65B is substantially below that of other unfinetuned models, suggesting that finetuning on the OASST1 dataset effectively reduces biases present in the original LLaMA model. Despite these promising findings, the extent to which {Guanaco} mitigates other bias categories remains uncertain. A broader, in-depth bias analysis for {Guanaco} and comparable conversational agents is left as future work.

Another shortcoming concerns the absence of experiments involving alternative bit-precisions—such as using 3-bit base models—or a range of adapter techniques. While LoRA was selected due to extensive evidence of its robustness, many Parameter Efficient FineTuning (PEFT) approaches have also demonstrated effectiveness, although their scalability to very large models is undetermined. Alternative adapters may outperform LoRA under certain conditions. Furthermore, given that post-quantization finetuning can recover most quantization-induced losses, even more aggressive compression (e.g., 3-bit GPTQ quantization paired with LoRA) could potentially match the results of full 16-bit finetuning once finetuned.

\section{Broader Impacts}

Our \textsc{QLoRA} approach is the first to facilitate the finetuning of 33B parameter models on a single consumer GPU and 65B parameter models on a single professional GPU, all while maintaining performance on par with full finetuning baselines. The results show that our top-performing 33B model, trained with the Open Assistant dataset, achieves competitiveness with ChatGPT as measured by the Vicuna benchmark. As instruction-based finetuning is pivotal for adapting pretrained LLMs into conversational agents resembling ChatGPT, we anticipate that our technique will broaden the accessibility and prevalence of finetuning, especially for researchers with limited computational resources—an important step forward for democratizing advanced NLP tools. Ultimately, \textsc{QLoRA} offers a means to reduce disparities in resources between major companies and smaller teams working with standard consumer-grade hardware.

An additional avenue of influence lies in the deployment of models onto mobile devices. We propose that our \textsc{QLoRA} technique could represent a crucial advancement, making it feasible to finetune LLMs directly on smartphones and other resource-limited hardware. Although it has previously been demonstrated that 7B-parameter models can be executed on mobile phones, \textsc{QLoRA} uniquely enables these models to be finetuned in such contexts. Our projections suggest that, on an iPhone 12 Plus, \textsc{QLoRA} would facilitate the finetuning of approximately 3 million tokens overnight while the device is plugged in. Even though finetuned 7B models may not yet match the performance levels of ChatGPT, we contend that their current quality is sufficient to unlock innovative use cases previously constrained by limitations in privacy or LLM effectiveness. In particular, \textsc{QLoRA} supports privacy-conscious LLM applications, empowering individuals to control their own models and data while also simplifying LLM deployment.

Nonetheless, it is important to recognize that finetuning is a technology with dual-use potential, capable of both beneficial and malicious applications. The widespread adoption of LLMs presents various risks \citep{bommasani2021opportunities,bender2021dangers}, yet we argue that democratizing access to this technology—which is rapidly gaining global prevalence—will promote broader and more autonomous evaluation, as opposed to consolidating LLM capabilities exclusively within major technology firms that withhold models and code from public scrutiny.

In summary, we anticipate that \textsc{QLoRA} will generate significant positive effects by substantially lowering the barriers to finetuning high-quality LLMs, thus promoting broader and simpler accessibility.